% Shared preamble for all subjects.
%
% Each subject's main.tex starts with:
%   % Shared preamble for all subjects.
%
% Each subject's main.tex starts with:
%   % Shared preamble for all subjects.
%
% Each subject's main.tex starts with:
%   % Shared preamble for all subjects.
%
% Each subject's main.tex starts with:
%   \input{../shared/preamble}
% followed by any packages/commands specific to that subject (e.g. tikz
% libraries, circuitikz, \dbar), then \begin{document}.

% \documentclass[openany,oneside]{book}

\usepackage{amsmath, amsthm, amssymb, amsfonts}
\usepackage{mathtools}
\usepackage{booktabs}
\usepackage{thmtools}
\usepackage{graphicx}
\usepackage{setspace}
\usepackage{geometry}
\usepackage{float}
\usepackage{hyperref}
\usepackage[utf8]{inputenc}
\usepackage[english]{babel}
\usepackage{framed}
\usepackage[dvipsnames]{xcolor}
\usepackage{tcolorbox}
\usepackage{bm}
\usepackage{physics}
\usepackage{enumitem}
\usepackage{cancel}
\usepackage{titlesec}

\colorlet{LightGray}{White!90!Periwinkle}
\colorlet{LightOrange}{Orange!15}
\colorlet{LightGreen}{Green!15}

\newcommand{\HRule}[1]{\rule{\linewidth}{#1}}

% breakable lets theorem/definition/formula/fact/example boxes split
% across a page boundary instead of being pushed whole onto the next
% page (a major source of the blank vertical space fixed below).
\tcbset{breakable, before skip=8pt, after skip=8pt}

\declaretheoremstyle[name=Theorem,]{thmsty}
\declaretheorem[style=thmsty,numberwithin=section]{theorem}
\tcolorboxenvironment{theorem}{colback=LightGray}

\declaretheoremstyle[name=Definition,]{defsty}
\declaretheorem[style=defsty,numberlike=theorem]{definition}
\tcolorboxenvironment{definition}{colback=LightGray}

\declaretheoremstyle[name=Formula,]{formsty}
\declaretheorem[style=formsty,numberlike=theorem]{formula}
\tcolorboxenvironment{formula}{colback=LightGray}

\declaretheoremstyle[name=Important Fact,]{factsty}
\declaretheorem[style=factsty,numberlike=theorem]{fact}
\tcolorboxenvironment{fact}{colback=LightGray}

\colorlet{ExampleGray}{black!8}
\declaretheoremstyle[name=Example,]{examsty}
\declaretheorem[style=examsty,numberlike=theorem]{example}
\tcolorboxenvironment{example}{colback=ExampleGray, colframe=ExampleGray, arc=6pt, boxrule=0pt}

\titleformat{\chapter}[display]
  {\normalfont\huge\bfseries}
  {\chaptertitlename\ \thechapter}
  {12pt}
  {\Huge}
\titlespacing*{\chapter}{0pt}{20pt}{16pt}
\titlespacing*{\section}{0pt}{10pt}{4pt}
\titlespacing*{\subsection}{0pt}{8pt}{3pt}
\titlespacing*{\subsubsection}{0pt}{6pt}{2pt}

\setlist{noitemsep, topsep=4pt, parsep=2pt}

\setstretch{1.2}
\geometry{
    textheight=9in,
    textwidth=6.5in,
    top=1in,
    headheight=12pt,
    headsep=25pt,
    footskip=30pt
}

% book defaults to \flushbottom, which stretches interline glue to make
% every page exactly textheight tall; with tcolorbox environments using
% fixed (non-stretchable) before/after skips, that stretch concentrates
% into large uneven gaps between ordinary paragraphs. Let pages end
% wherever the content ends instead.
\raggedbottom

% Without this, amsmath refuses to break a multi-line align/gather
% across a page boundary, so a long derivation that doesn't quite fit
% gets pushed whole onto the next page, leaving a large blank remainder
% at the bottom of the current one.
\allowdisplaybreaks

\newcommand{\ihat}{\boldsymbol{\hat{\textbf{i}}}}
\newcommand{\jhat}{\boldsymbol{\hat{\textbf{j}}}}
\newcommand{\khat}{\boldsymbol{\hat{\textbf{k}}}}

\newcommand{\rhat}{\boldsymbol{\hat{\textbf{r}}}}
\newcommand{\tthat}{\hat{\boldsymbol{\theta}}}

\newcommand{\vect}[1]{\boldsymbol{\overrightarrow{#1}}}
\newcommand{\svect}[1]{\boldsymbol{\vec{#1}}}
\newcommand{\uvect}[1]{\boldsymbol{\hat{\textbf{#1}}}}

\newcommand{\defeq}{\vcentcolon=}

% followed by any packages/commands specific to that subject (e.g. tikz
% libraries, circuitikz, \dbar), then \begin{document}.

% \documentclass[openany,oneside]{book}

\usepackage{amsmath, amsthm, amssymb, amsfonts}
\usepackage{mathtools}
\usepackage{booktabs}
\usepackage{thmtools}
\usepackage{graphicx}
\usepackage{setspace}
\usepackage{geometry}
\usepackage{float}
\usepackage{hyperref}
\usepackage[utf8]{inputenc}
\usepackage[english]{babel}
\usepackage{framed}
\usepackage[dvipsnames]{xcolor}
\usepackage{tcolorbox}
\usepackage{bm}
\usepackage{physics}
\usepackage{enumitem}
\usepackage{cancel}
\usepackage{titlesec}

\colorlet{LightGray}{White!90!Periwinkle}
\colorlet{LightOrange}{Orange!15}
\colorlet{LightGreen}{Green!15}

\newcommand{\HRule}[1]{\rule{\linewidth}{#1}}

% breakable lets theorem/definition/formula/fact/example boxes split
% across a page boundary instead of being pushed whole onto the next
% page (a major source of the blank vertical space fixed below).
\tcbset{breakable, before skip=8pt, after skip=8pt}

\declaretheoremstyle[name=Theorem,]{thmsty}
\declaretheorem[style=thmsty,numberwithin=section]{theorem}
\tcolorboxenvironment{theorem}{colback=LightGray}

\declaretheoremstyle[name=Definition,]{defsty}
\declaretheorem[style=defsty,numberlike=theorem]{definition}
\tcolorboxenvironment{definition}{colback=LightGray}

\declaretheoremstyle[name=Formula,]{formsty}
\declaretheorem[style=formsty,numberlike=theorem]{formula}
\tcolorboxenvironment{formula}{colback=LightGray}

\declaretheoremstyle[name=Important Fact,]{factsty}
\declaretheorem[style=factsty,numberlike=theorem]{fact}
\tcolorboxenvironment{fact}{colback=LightGray}

\colorlet{ExampleGray}{black!8}
\declaretheoremstyle[name=Example,]{examsty}
\declaretheorem[style=examsty,numberlike=theorem]{example}
\tcolorboxenvironment{example}{colback=ExampleGray, colframe=ExampleGray, arc=6pt, boxrule=0pt}

\titleformat{\chapter}[display]
  {\normalfont\huge\bfseries}
  {\chaptertitlename\ \thechapter}
  {12pt}
  {\Huge}
\titlespacing*{\chapter}{0pt}{20pt}{16pt}
\titlespacing*{\section}{0pt}{10pt}{4pt}
\titlespacing*{\subsection}{0pt}{8pt}{3pt}
\titlespacing*{\subsubsection}{0pt}{6pt}{2pt}

\setlist{noitemsep, topsep=4pt, parsep=2pt}

\setstretch{1.2}
\geometry{
    textheight=9in,
    textwidth=6.5in,
    top=1in,
    headheight=12pt,
    headsep=25pt,
    footskip=30pt
}

% book defaults to \flushbottom, which stretches interline glue to make
% every page exactly textheight tall; with tcolorbox environments using
% fixed (non-stretchable) before/after skips, that stretch concentrates
% into large uneven gaps between ordinary paragraphs. Let pages end
% wherever the content ends instead.
\raggedbottom

% Without this, amsmath refuses to break a multi-line align/gather
% across a page boundary, so a long derivation that doesn't quite fit
% gets pushed whole onto the next page, leaving a large blank remainder
% at the bottom of the current one.
\allowdisplaybreaks

\newcommand{\ihat}{\boldsymbol{\hat{\textbf{i}}}}
\newcommand{\jhat}{\boldsymbol{\hat{\textbf{j}}}}
\newcommand{\khat}{\boldsymbol{\hat{\textbf{k}}}}

\newcommand{\rhat}{\boldsymbol{\hat{\textbf{r}}}}
\newcommand{\tthat}{\hat{\boldsymbol{\theta}}}

\newcommand{\vect}[1]{\boldsymbol{\overrightarrow{#1}}}
\newcommand{\svect}[1]{\boldsymbol{\vec{#1}}}
\newcommand{\uvect}[1]{\boldsymbol{\hat{\textbf{#1}}}}

\newcommand{\defeq}{\vcentcolon=}

% followed by any packages/commands specific to that subject (e.g. tikz
% libraries, circuitikz, \dbar), then \begin{document}.

% \documentclass[openany,oneside]{book}

\usepackage{amsmath, amsthm, amssymb, amsfonts}
\usepackage{mathtools}
\usepackage{booktabs}
\usepackage{thmtools}
\usepackage{graphicx}
\usepackage{setspace}
\usepackage{geometry}
\usepackage{float}
\usepackage{hyperref}
\usepackage[utf8]{inputenc}
\usepackage[english]{babel}
\usepackage{framed}
\usepackage[dvipsnames]{xcolor}
\usepackage{tcolorbox}
\usepackage{bm}
\usepackage{physics}
\usepackage{enumitem}
\usepackage{cancel}
\usepackage{titlesec}

\colorlet{LightGray}{White!90!Periwinkle}
\colorlet{LightOrange}{Orange!15}
\colorlet{LightGreen}{Green!15}

\newcommand{\HRule}[1]{\rule{\linewidth}{#1}}

% breakable lets theorem/definition/formula/fact/example boxes split
% across a page boundary instead of being pushed whole onto the next
% page (a major source of the blank vertical space fixed below).
\tcbset{breakable, before skip=8pt, after skip=8pt}

\declaretheoremstyle[name=Theorem,]{thmsty}
\declaretheorem[style=thmsty,numberwithin=section]{theorem}
\tcolorboxenvironment{theorem}{colback=LightGray}

\declaretheoremstyle[name=Definition,]{defsty}
\declaretheorem[style=defsty,numberlike=theorem]{definition}
\tcolorboxenvironment{definition}{colback=LightGray}

\declaretheoremstyle[name=Formula,]{formsty}
\declaretheorem[style=formsty,numberlike=theorem]{formula}
\tcolorboxenvironment{formula}{colback=LightGray}

\declaretheoremstyle[name=Important Fact,]{factsty}
\declaretheorem[style=factsty,numberlike=theorem]{fact}
\tcolorboxenvironment{fact}{colback=LightGray}

\colorlet{ExampleGray}{black!8}
\declaretheoremstyle[name=Example,]{examsty}
\declaretheorem[style=examsty,numberlike=theorem]{example}
\tcolorboxenvironment{example}{colback=ExampleGray, colframe=ExampleGray, arc=6pt, boxrule=0pt}

\titleformat{\chapter}[display]
  {\normalfont\huge\bfseries}
  {\chaptertitlename\ \thechapter}
  {12pt}
  {\Huge}
\titlespacing*{\chapter}{0pt}{20pt}{16pt}
\titlespacing*{\section}{0pt}{10pt}{4pt}
\titlespacing*{\subsection}{0pt}{8pt}{3pt}
\titlespacing*{\subsubsection}{0pt}{6pt}{2pt}

\setlist{noitemsep, topsep=4pt, parsep=2pt}

\setstretch{1.2}
\geometry{
    textheight=9in,
    textwidth=6.5in,
    top=1in,
    headheight=12pt,
    headsep=25pt,
    footskip=30pt
}

% book defaults to \flushbottom, which stretches interline glue to make
% every page exactly textheight tall; with tcolorbox environments using
% fixed (non-stretchable) before/after skips, that stretch concentrates
% into large uneven gaps between ordinary paragraphs. Let pages end
% wherever the content ends instead.
\raggedbottom

% Without this, amsmath refuses to break a multi-line align/gather
% across a page boundary, so a long derivation that doesn't quite fit
% gets pushed whole onto the next page, leaving a large blank remainder
% at the bottom of the current one.
\allowdisplaybreaks

\newcommand{\ihat}{\boldsymbol{\hat{\textbf{i}}}}
\newcommand{\jhat}{\boldsymbol{\hat{\textbf{j}}}}
\newcommand{\khat}{\boldsymbol{\hat{\textbf{k}}}}

\newcommand{\rhat}{\boldsymbol{\hat{\textbf{r}}}}
\newcommand{\tthat}{\hat{\boldsymbol{\theta}}}

\newcommand{\vect}[1]{\boldsymbol{\overrightarrow{#1}}}
\newcommand{\svect}[1]{\boldsymbol{\vec{#1}}}
\newcommand{\uvect}[1]{\boldsymbol{\hat{\textbf{#1}}}}

\newcommand{\defeq}{\vcentcolon=}

% followed by any packages/commands specific to that subject (e.g. tikz
% libraries, circuitikz, \dbar), then \begin{document}.

% \documentclass[openany,oneside]{book}

\usepackage{amsmath, amsthm, amssymb, amsfonts}
\usepackage{mathtools}
\usepackage{booktabs}
\usepackage{thmtools}
\usepackage{graphicx}
\usepackage{setspace}
\usepackage{geometry}
\usepackage{float}
\usepackage{hyperref}
\usepackage[utf8]{inputenc}
\usepackage[english]{babel}
\usepackage{framed}
\usepackage[dvipsnames]{xcolor}
\usepackage{tcolorbox}
\usepackage{bm}
\usepackage{physics}
\usepackage{enumitem}
\usepackage{cancel}
\usepackage{titlesec}

\colorlet{LightGray}{White!90!Periwinkle}
\colorlet{LightOrange}{Orange!15}
\colorlet{LightGreen}{Green!15}

\newcommand{\HRule}[1]{\rule{\linewidth}{#1}}

% breakable lets theorem/definition/formula/fact/example boxes split
% across a page boundary instead of being pushed whole onto the next
% page (a major source of the blank vertical space fixed below).
\tcbset{breakable, before skip=8pt, after skip=8pt}

\declaretheoremstyle[name=Theorem,]{thmsty}
\declaretheorem[style=thmsty,numberwithin=section]{theorem}
\tcolorboxenvironment{theorem}{colback=LightGray}

\declaretheoremstyle[name=Definition,]{defsty}
\declaretheorem[style=defsty,numberlike=theorem]{definition}
\tcolorboxenvironment{definition}{colback=LightGray}

\declaretheoremstyle[name=Formula,]{formsty}
\declaretheorem[style=formsty,numberlike=theorem]{formula}
\tcolorboxenvironment{formula}{colback=LightGray}

\declaretheoremstyle[name=Important Fact,]{factsty}
\declaretheorem[style=factsty,numberlike=theorem]{fact}
\tcolorboxenvironment{fact}{colback=LightGray}

\colorlet{ExampleGray}{black!8}
\declaretheoremstyle[name=Example,]{examsty}
\declaretheorem[style=examsty,numberlike=theorem]{example}
\tcolorboxenvironment{example}{colback=ExampleGray, colframe=ExampleGray, arc=6pt, boxrule=0pt}

\titleformat{\chapter}[display]
  {\normalfont\huge\bfseries}
  {\chaptertitlename\ \thechapter}
  {12pt}
  {\Huge}
\titlespacing*{\chapter}{0pt}{20pt}{16pt}
\titlespacing*{\section}{0pt}{10pt}{4pt}
\titlespacing*{\subsection}{0pt}{8pt}{3pt}
\titlespacing*{\subsubsection}{0pt}{6pt}{2pt}

\setlist{noitemsep, topsep=4pt, parsep=2pt}

\setstretch{1.2}
\geometry{
    textheight=9in,
    textwidth=6.5in,
    top=1in,
    headheight=12pt,
    headsep=25pt,
    footskip=30pt
}

% book defaults to \flushbottom, which stretches interline glue to make
% every page exactly textheight tall; with tcolorbox environments using
% fixed (non-stretchable) before/after skips, that stretch concentrates
% into large uneven gaps between ordinary paragraphs. Let pages end
% wherever the content ends instead.
\raggedbottom

% Without this, amsmath refuses to break a multi-line align/gather
% across a page boundary, so a long derivation that doesn't quite fit
% gets pushed whole onto the next page, leaving a large blank remainder
% at the bottom of the current one.
\allowdisplaybreaks

\newcommand{\ihat}{\boldsymbol{\hat{\textbf{i}}}}
\newcommand{\jhat}{\boldsymbol{\hat{\textbf{j}}}}
\newcommand{\khat}{\boldsymbol{\hat{\textbf{k}}}}

\newcommand{\rhat}{\boldsymbol{\hat{\textbf{r}}}}
\newcommand{\tthat}{\hat{\boldsymbol{\theta}}}

\newcommand{\vect}[1]{\boldsymbol{\overrightarrow{#1}}}
\newcommand{\svect}[1]{\boldsymbol{\vec{#1}}}
\newcommand{\uvect}[1]{\boldsymbol{\hat{\textbf{#1}}}}

\newcommand{\defeq}{\vcentcolon=}
